% =====================================================================
% 1.1 Motivation and research question
% No displayed equations here: the model is stated formally in Section 3.
% =====================================================================
\subsection{MOTIVATION AND RESEARCH QUESTION}
\label{subsec:motivation}

A single default rarely threatens a lender; many defaults occurring together do. Corporate defaults cluster because the asset values of different firms respond to the same macroeconomic and market shocks. The loss distribution of a credit portfolio is therefore strongly skewed: most of its mass lies close to zero, while its right tail is long and is governed far more by the \emph{dependence} among obligors than by their individual default probabilities \citep[Ch.~8]{McNeil05}. Risk measures evaluated deep in that tail, such as value-at-risk (VaR) and expected shortfall (ES) at the $99.9\%$ level, are the quantities most exposed to how this dependence is modelled.

This sensitivity has direct regulatory consequences. The Basel internal ratings-based (IRB) approach derives credit capital from the $99.9\%$ quantile of the loss of a portfolio driven by a single Gaussian systematic factor \citep[Sect.~6.4]{DeGiorgiFilkova26}, and widely used industry models such as CreditMetrics are built on Gaussian factors as well \citep[Sect.~8.2]{McNeil05}. Gaussian dependence, however, is asymptotically tail-independent: whatever the correlation below one, the probability that two variables take extreme values \emph{together} vanishes in the limit \citep[Sect.~5.2.3]{McNeil05}. Equity markets tend to behave in the opposite way, falling together in crises such as autumn 2008 and March 2020; Section~\ref{sec:data} documents this pattern in our own data. A model that rules out joint extremes may thus understate capital in exactly the states of the world in which capital is needed.

Copulas offer a clean way to address this issue. By Sklar's theorem, every joint distribution can be split into its marginal distributions, which describe how risky each variable is on its own, and a copula, which describes how the variables move together, including in the tails \citep[Ch.~5]{McNeil05}. In credit risk this separation is particularly useful: once the individual default probabilities are fixed, the copula of the firms' asset returns alone determines how defaults cluster \citep[Sect.~8.3]{McNeil05}. Its impact can be large. In a simulation study by \citet[Table~8.4]{McNeil05}, replacing a Gaussian copula by a $t$-copula with $\nu=10$ degrees of freedom, with default probabilities and asset correlations unchanged, raises the $99\%$ quantile of the number of defaults by a factor of almost six.

We investigate this question for a portfolio of $100$ corporate bonds issued by firms with branches in Switzerland and the United States, over a horizon of one week. Default is modelled with a threshold model in the spirit of Merton's structural model \citep[Ch.~2]{DeGiorgi26}: a firm defaults when its asset-value return falls below a critical threshold (Appendix~A). Each firm's asset return depends on two observable systematic factors, the weekly returns of the Swiss Performance Index (SPI) and of the S\&P~500 (SPX), and on an independent firm-specific shock \citep[Sect.~3]{DeGiorgiFilkova26}. Because the firm-specific shocks are independent, the two equity indices are the only channel through which defaults can cluster, and the copula question moves from the firms to the two market factors. This also sets our model apart from the experiment above: dependence beyond the Gaussian case enters only through the systematic component of each firm, while the Gaussian firm-specific shock dilutes it. Whether a substantial tail effect survives this dilution is an empirical question.

The central question of this paper is therefore:
\begin{quote}
\itshape Holding every counterparty's default probability fixed, how much of the $99.9\%$ VaR and ES of the portfolio is driven by the joint distribution assumed for the Swiss and US equity factors, and in particular by their dependence in the tails?
\end{quote}
